% Chapter 3, section 5: Workflows: one step, one commit
% Write in plain LaTeX. Code of a guide step: \guide{<step>}; commands: \begin{shell}...\end{shell};
% a screenshot: \screenshot{file}{caption}; a gap in the draft: \placeholder{...}.
% See book/latex/queen.sty for the full list.
%
% Draft by Claude, 2026-10-24. The code is cut from examples/tutorials/js/book-use-cases.mjs, which
% checks every result quoted here against a 2.6.0 node.

\section{Workflows: one step, one commit}

An order moves through states: paid, shipping, shipped. A redelivered event, two workers, or an
event that arrives late must never move it twice or backwards, and a step must never be half done.

The order's state lives in KV. Each step is one transaction: it acks the event, writes the new
state only if the state is still the version the worker read (\texttt{expect}), and pushes the next
event. If the version moved, the whole transaction rolls back, and nothing is written.

\queenfig{order-workflow}{A state machine whose steps cannot race.}

\guide{uc-order}

The step moved \texttt{order-2001} to \texttt{shipping} and pushed one shipping event. The same
step sent again from the old version rolled back: still one event.

\why{the ack, the new state and the next event commit together or not at all. That replaces the
outbox table, the idempotency table and the lock you would put around them.}

One thing to know: \texttt{commit()} returns \texttt{success: false} for a rollback, it does not
throw. Read it.
